\section{Rethinking the filters: global metrics bias toward rigid folds}\label{sec:filtering}
With the pipeline reproduced and trusted (Section~\ref{sec:methods}), the first thing that
needed rethinking was the filtering. The original filters score a design with two global
continuous metrics, overall RMSD and mean PAE, that treat the epitope and the scaffold as
one unit. This is a reasonable inheritance from standard de novo practice rather than an
arbitrary choice, but it biases selection toward globally rigid, well-ordered folds and
silently rejects designs whose scaffold is flexible while the epitope is well positioned
(or the reverse). Decomposing each metric into epitope- and scaffold-specific terms
(Section~\ref{sec:decomp}) is what lets us see, and eventually recover, that
rejected-but-viable space.

To compare design sets on a common footing we use epitope-chunk RMSD $< \SI{2.5}{\angstrom}$
as a single bar. All three tiled 12-mer antigens pass it more often than the pooled DP3 set:

\begin{center}
\begin{tabular}{lcc}
\toprule
set & pass rate ($<\SI{2.5}{\angstrom}$) & median epitope RMSD (\si{\angstrom}) \\
\midrule
DP3 / mAb (pooled) & 30.0\% & 5.00 \\
6M0J               & 47.0\% & 2.65 \\
1D2K               & 60.2\% & 1.85 \\
4WAT               & 84.4\% & 0.43 \\
\bottomrule
\end{tabular}
\end{center}

\noindent
Even the hardest 12-mer antigen (6M0J) clears the bar more often than DP3. The full
six-metric distribution comparison is Fig.~\ref{fig:passrates}.

\begin{figure}[h]\centering
\figorbox{dp3_vs_12mer_ungated.png}{0.95}
\caption{Ungated metric distributions (full design sets), DP3 mAb versus the three 12-mer
antigens, across epitope RMSD, global mean PAE (the four-filter metric), epitope PAE (the
decomposed metric the composite scores on), overall RMSD, pTM, AF3 clashing residues
(antibody set only), and cylinder clashes. Dashed line marks the median; red marks the
\SI{2.5}{\angstrom} epitope-RMSD gate with the per-set pass rate. Regenerated by
\texttt{notebooks/dp3\_vs\_12mer\_distributions.ipynb}.}
\label{fig:passrates}
\end{figure}

\paragraph{Caveats (important for interpretation).}
This comparison is not yet apples-to-apples. The 12-mer epitope is a single contiguous
12-residue stretch, whereas DP3 epitopes are larger and often multi-island, so 12-mer
RMSD/PAE are mechanically lower in part by definition, and the pooled DP3 30\% also averages
over heterogeneous epitopes. The clean control is to stratify within DP3 by island count
and size (Section~\ref{sec:open}). Cylinder counts are likewise not comparable across sets as
absolute numbers. The 12-mer cylinder medians ($\sim$31--32) exceed DP3 ($\sim$21) largely
because the designs are a fixed length ($\sim$104 residues) while the 12-mer epitope is small
(12 residues): that leaves $\sim$91 non-epitope scaffold residues whose C$\alpha$ atoms can
fall in the fixed cylinder, versus fewer for DP3's larger epitopes. It is a scaffold-mass
artifact rather than more genuine clashing, which is why the score applies the native-aware
carve and ranks the cylinder per antigen rather than against an absolute cutoff. As a sanity
check, the DP3 cylinder median ($\approx 21$) tracks its real \texttt{af3\_n\_clash\_res}
median ($\approx 19$).

\paragraph{Global versus epitope PAE.}
The PAE column is a clean illustration of why decomposition matters, and the two are easy to
confuse. The four-filter (and Lawson's original) uses the \emph{global} mean PAE
($<5$); on this DP3 set that distribution sits high, median $13.2$ with only $\sim$7\% of
designs below 5. That global value is dragged up by the \emph{scaffold} region, which
AlphaFold3 predicts less confidently. The epitope region itself is much better resolved
(epitope PAE median $9.1$, $\sim$23\% below 5). A design can therefore hold its epitope well
and still fail the global PAE filter on scaffold uncertainty that may not matter for binding,
which is precisely the bias the decomposed metrics are meant to relax. For this reason the
figure shows both columns, and the composite scores the epitope-specific PAE rather than the
global one.

\section{Composite scoring and per-epitope selection}\label{sec:composite}
Decomposed metrics give more knobs, which raised the next question: how to combine them into
a single ranking so we can take the best designs per epitope. The answer is the
config-driven composite scorer (Section~\ref{sec:methods}). It gates on epitope-chunk RMSD,
transforms each metric to a higher-is-better score within its scope (per antigen for the
12-mer set, pooled for DP3), weight-sums, and keeps the top designs per epitope. Run over
the full 12-mer set, the gate keeps \num{54539} of \num{82712} designs ($\approx 65.9\%$
pooled) before ranking \texttt{[regen]}, consistent with the per-antigen rates above.

The scorer is deliberately simple: one metric per input, each with its own activation,
weighted and summed, which makes it a single-layer perceptron with hand-set weights. It is
still in flux. The per-metric transform is currently a percentile within the population, so
a design's score depends on the cohort it is ranked against. Whether to switch to a
population-independent sigmoid (an absolute, saturating per-metric score) is an open choice
carried in Section~\ref{sec:open}.

\section{Scaffolding without an antibody: the cylinder surrogate}\label{sec:noab}
Most targets of interest have no solved antibody, so the four-filter ground truth does not
apply, because there is no antibody to clash against. We extend the workflow to these
antigens by tiling them into 12-residue epitopes and replacing the antibody-clash filter
with a native-antigen-aware cylinder surrogate (Section~\ref{sec:methods}): a cylinder
fitted to the epitope plane along the antibody-approach normal that counts non-epitope
scaffold mass intruding into the volume an approaching antibody would need
(Fig.~\ref{fig:cylinder}).

\begin{figure}[h]\centering
\figorbox{epitope_cylinder.png}{0.5}
\caption{The accessibility cylinder. Epitope residues (red) sit on the scaffold (magenta);
the epitope C$\alpha$ atoms define a plane (orange grid) whose outward normal (blue arrow,
the antibody-approach direction) orients the cylinder (cyan). Scaffold C$\alpha$ atoms
inside the cylinder count against accessibility. Rendered in VMD by
\texttt{known\_antigen/analysis/cylinder\_vis/visualize\_cylinder.tcl}.}
\label{fig:cylinder}
\end{figure}

A naive cylinder over-penalizes, because some scaffold mass sits exactly where the native
antigen itself already sits, a volume a real antibody clearly tolerates since it binds the
native antigen there. The native-aware carve fixes this: among truly clash-free designs, we
exclude scaffold C$\alpha$ atoms within \SI{1}{\angstrom} of a native-antigen heavy atom.
Concretely, of the \num{3186} DP3 designs with zero true antibody clash, the number the
plain cylinder flags as heavily penalized (count $\ge 10$) falls from \num{1112} to
\num{384} once the native footprint is carved out, which is \textbf{\num{728} false-positive
penalties removed}, while genuinely clashing designs are left essentially untouched
(Figs.~\ref{fig:fp} and~\ref{fig:fpsel}).

\begin{figure}[h]\centering
\figorbox{fp_reduction.png}{0.8}
\caption{Cylinder count among clash-free DP3 designs (zero AlphaFold3 antibody clash,
$n=\num{3186}$), before (grey) and after (red) subtracting the native-antigen volume. The
heavily penalized tail (count $\ge 10$) collapses from \num{1112} to \num{384}. Regenerated
by \texttt{episcaf\_analysis/viz/plot\_fp\_reduction.py --native metrics\_native\_cyl.csv
--mode dist}.}
\label{fig:fp}
\end{figure}

\begin{figure}[h]\centering
\figorbox{fp_selectivity.png}{0.8}
\caption{The carve removes flags selectively. Fraction of designs flagged (cylinder
$\ge 10$) across bins of true antibody clash, before versus after the carve. The gap is
large where there is no real clash (false positives removed) and closes where the clash is
real (true signal retained). \texttt{--mode selectivity}.}
\label{fig:fpsel}
\end{figure}

\paragraph{A pipeline constraint worth recording: crystal gaps.}
Epitope fidelity for the 12-mer set is measured by aligning the AlphaFold3 epitope onto the
native crystal epitope, matched by residue number (\texttt{build\_12mer\_metrics.py}). When
a tile's epitope residues are unresolved in the crystal, meaning a gap in the deposited
structure, the native C$\alpha$ anchors are missing, the Kabsch alignment cannot be formed,
and the tile is dropped with status \texttt{crystal\_epi\_incomplete}. This is a real
constraint on which 12-mer epitopes are scorable at all: an antigen whose crystal has gaps
over the tiled region contributes fewer (or no) scorable tiles, independent of design
quality. The mechanism is confirmed in code; the specific antigen(s) and tile counts
affected are not in the local ok-only export and remain to be pulled from the cluster
pre-filter status counts \texttt{[verify]}.

\subsection{Interpreting the cylinder: where passing designs land}
The cylinder is easiest to read by overlaying, for each metric, the full design population
against the designs we would actually advance (Fig.~\ref{fig:overlay}). For DP3 the advanced
set is the four-filter passes (\num{751} designs); for DP4 it is the top three per epitope by
composite score. On every geometry and confidence metric the passing designs concentrate
where expected: low RMSD, low PAE, and high pTM. The cylinder column is the informative one.
The passing DP3 designs, which we know the antibody accepts, sit at low cylinder counts
(median near 10), whereas the DP4 top-3 sit higher (medians near 25--30). This is the same
scaffold-mass effect noted earlier rather than a sign that the DP4 designs are worse: with a
12-residue epitope on a $\sim$104-residue design, more non-epitope scaffold mass can fall in
the fixed cylinder even among the designs we would pick.
It is a direct argument for leaning on the native-aware carve and per-antigen scoring rather
than an absolute cylinder cutoff shared across sets, and the comparison gives us a concrete
reference for what a tolerable cylinder count looks like (the DP3 passing range).

The cylinder earns this role. On DP3, where the true antibody clash is known, the cylinder
predicts clash-free designs (\texttt{af3\_n\_clash\_res}=0) with AUC $0.93$
(\texttt{cylinder\_native\_aware} $0.935$, slightly better than the plain count $0.924$;
clash-free vs clashing medians 5 vs 22). It is therefore a strong stand-in for the filter we
lose on DP4, and the native-aware carve helps the ranking, not just the optics
(\texttt{episcaf\_analysis/cylinder\_surrogate\_auc.py}).

\begin{figure}[h]\centering
\figorbox{passes_overlay.png}{1.0}
\caption{All designs (grey, left axis) and passing designs (color, right colored axis), both
in true counts, per metric and per row; the dashed line marks the median of the passes. The
twin axes keep the full population and the much smaller passing set each at its own scale, so
the shape of the accepted distribution and where it lands are both visible. DP3 passes are
the four-filter set; DP4 passes are the top 3 per epitope by composite score. Regenerated by
\texttt{episcaf\_analysis/viz/plot\_passes\_overlay.py}.}
\label{fig:overlay}
\end{figure}
